\chapter{人工智慧的使用}
\label{app:ai-usage}

\section{開發過程中的生成式 AI}

\textbf{使用的工具。}主要使用 Claude Code（Claude Opus 4.8、5.5），讓模型直接在專案目錄中讀寫程式與文件；另以 ChatGPT（Terra 5.6）作為不同來源的第二個模型，用於討論與交叉檢查。

\textbf{程式。}系統的程式碼幾乎全由 AI 撰寫。組員負責決定架構與設計，包括模組切分、階段之間的資料契約、採用哪些既有模型、實驗的評估方式，以及介面的互動與視覺設計，實作則交給 AI。實際流程通常是先請 AI 提出方案而不動手，由組員取捨後再實作，完成後由 AI 列出需要人工確認的項目，組員實際操作系統檢查，確認無誤才提交版本。

\textbf{研究。}AI 用於搜尋文獻與導讀論文，例如整理 BST 在 ShuttleSet 上的資料切分方式，以確認哪些比賽未參與其訓練。也用於部分演算法的優化、以小型實驗驗證想法是否可行，以及把中間結果視覺化，例如將骨架、軌跡與擊球時刻疊回影片逐幀檢視。

\textbf{報告與展示。}AI 參與了討論章節的內容整理、全文順稿與排版，並撰寫報告圖片的生成程式。專題競賽的影片分鏡、模組介紹動畫的程式與海報版面規劃也由 AI 協助。報告中的數據皆來自實驗程式的輸出，並非由 AI 生成。

\section{確保正確性的做法}

AI 產出的程式與文字可能看似合理卻有錯，因此我們依影響程度採取不同的檢查方式：

\begin{enumerate}
  \item \textbf{評分程式由人閱讀。}實驗的評分程式決定報告中所有準確率數字，因此要求 AI 以最直白的寫法撰寫，並由組員自行讀過確認邏輯。其他必要的程式也同樣自行閱讀。
  \item \textbf{以不同模型交叉檢查。}重要的程式與文件會交由另一個模型審查。例如第二個模型在審查分段實驗時發現，標註工具輸出的 1354 筆真值，其秒數欄位全部少算一幀。
  \item \textbf{視覺化檢查。}不只看程式回報的數字，而是把結果疊回影片或畫成圖，以人眼確認是否合理，許多錯誤是在這一步才發現的。
  \item \textbf{標出依據。}要求 AI 在提出結論時附上依據，例如程式的檔案與行號、實驗輸出的檔案或原始論文的段落，方便回頭查證。
  \item \textbf{重點驗證結果。}關鍵數字重新執行評分程式確認，修改模組後也以既有的評估結果作為回歸檢查，確認沒有退步。
\end{enumerate}

AI 大幅縮短了實作所需的時間，讓小組能把心力放在設計與驗證上。代價是組員對部分實作細節並不熟悉，因此我們把人工閱讀集中在評分程式這類直接影響結論的部分。

\chapter{音訊實驗來源與待補資料}
\label{app:audio-sources}

模型、資料角色、指標來源、JSON hashes、實作對照與重現限制見 \file{experiments/audio_highlight/README.md}。待補內容包含 \texttt{match\_001} 至 \texttt{match\_004} 的真實賽事來源、完整 cheer／no-cheer 標註規範、多標註者一致性，以及 labels／features／結果的可重現封存。這些缺口不由比賽資料夾名稱或模型輸出推導。
